% Part: computability
% Chapter: computability-theory
% Section: def-functions-self-reference

\documentclass[../../../include/open-logic-section]{subfiles}

\begin{document}

\olfileid[fr]{cmp}{thy}{slf}
\olsection{Définir des fonctions par autoréférence}

Il est souvent utile de pouvoir définir une fonction
en termes d'elle-même. Par exemple, si $k$, $l$
et~$m$ sont des fonctions calculables totales, le théorème
du point fixe donne une fonction calculable
partielle~$f$ qui satisfait, pour tout~$y$,
\[
f(y) \simeq
\begin{cases}
k(y) & \text{si $l(y) = 0$} \\
f(m(y)) & \text{sinon.}
\end{cases}
\]
Plus précisément, on obtient~$f$ en posant
\[
g(x,y) \simeq
\begin{cases}
k(y) & \text{si $l(y) = 0$} \\
\cfind{x}(m(y)) & \text{sinon}
\end{cases}
\]
puis en appliquant le théorème du point fixe
pour trouver un indice~$e$ tel que
$\cfind{e}(y) \simeq g(e,y)$.

Pour un exemple concret, on peut définir la
fonction « plus grand commun diviseur »
$\fn{gcd}(u,v)$ par
\[
\fn{gcd}(u,v) \simeq
\begin{cases}
v & \text{si $u = 0$} \\
\fn{gcd}(\fn{mod}(v, u), u) & \text{sinon}
\end{cases}
\]
où $\fn{mod}(v,u)$ désigne le reste de la
division de~$v$ par~$u$. Le théorème du point
fixe montre que $\fn{gcd}$ est calculable
partielle. (On peut en effet ramener cet
exemple à la forme précédente en faisant coder
le couple $\tuple{u,v}$ par~$y$.) Une récurrence
sur~$u$ montre ensuite que cette fonction est,
en fait, totale : pour $u>0$, on a $0\leq\fn{mod}(v,u)<u$, et pour $u=0$ la valeur est immédiatement $v$ (donc $\fn{gcd}(0,0)=0$ selon cette convention).\footnote{Nous employons $\simeq$ dans l'application du théorème du point fixe, contrairement à l'égalité ordinaire de la source, pour couvrir les valeurs partielles avant la preuve de totalité.}

On peut bien sûr inventer des définitions
autoréférentielles beaucoup plus élaborées.
La plupart des langages de programmation
permettent, d'une manière ou d'une autre, de
définir une fonction en termes d'elle-même.
Il s'agit toutefois d'un pouvoir un peu moins
remarquable que celui de définir une fonction
à partir d'un \emph{indice} d'un algorithme qui
la calcule : c'est cette dernière possibilité
que donne le théorème du point fixe dans toute
sa généralité.

\end{document}
